\chapter{Turkish and Kurmanji in Germany and Turkey}\label{chap:key:2}
\section{Turkish as a heritage language in Germany}\label{sec:key:2.1}

Turkish is one of the largest immigrant communities in Germany, with nearly five million speakers of the Turkish language (\citealt{KaranfilŞavk2013}). Although migration from Turkey to Germany began approximately 60 years ago, and many Turkish families now reside in Germany for three or even four generations, the Turkish language continues to be regarded as a relatively vital language. Turkish speakers maintain their language due to several factors, including the presence of numerous religious and social institutions where Turkish is regularly spoken, easy access to Turkish media, ongoing contact with Turkey, and a relatively low rate of exogamous marriages \citep{Bhatia2004}.

In cities with large Turkish-speaking populations, such as Berlin, Cologne, and Frankfurt, people can use Turkish not only at home with family but also in informal public settings, including shops, cafes, and during leisure activities (\citealt{Hottmann2008}, \citealt{IefremenkoEtAl2021}). Furthermore, studies on language orientation within Turkish communities in Western Europe and Australia (\citealt{Yagmur2004}, \citealt{YagmurVanDeVijver2012}) have found that the Turkish community in Germany exhibits the highest levels of ethnic identification and the lowest levels of mainstream identification among the communities studied. This strong sense of ethnic identity contributes to the vitality of the Turkish language.

Immigration from Turkey to Germany began in the 1960s with a recruitment agreement between the two countries, allowing Turks to come to Germany as \textit{Gastarbeiter}, or `guest workers.' These individuals were primarily male laborers with little to no vocational training, originating from rural areas throughout Turkey \citep{Bayram2013}. Under the initial agreement, workers were expected to stay for only a short period before returning to their homeland. However, in the early 1970s, the law changed, permitting guest workers to remain in Germany and reunite with their families \citep{Angin2003}. This reunification policy led to a rapid increase in the size of the Turkish community in Germany and transformed the demographic composition from predominantly male workers to families with children \citep{Bayram2013}.

At the beginning of the guest workers' stay in Germany, they were regarded as temporary residents, leading to minimal efforts to facilitate their integration into German society. Furthermore, it was unclear whether schooling should be compulsory for children of guest workers \citep{Beck1999}. Subsequently, various educational models were adopted by different German states. For instance, Berlin placed foreign children, including those born in Germany, in \textit{Ausländerregelklassen} (`preparatory German language classes') due to their insufficient proficiency in German \citep{Pfaff1991}. In contrast, Bavaria and Baden-Württemberg implemented \textit{Muttersprachenunterricht} (`mother tongue education'), allowing children to study the history, arts, culture, geography, and religion of their parents’ country of origin \citep{Damanakis1983}. The primary aim of these classes was to prepare children for easier integration should they return to Turkey \citep{Damanakis1983}. However, both models ultimately failed, further hindering immigrant children’s integration into German society and their access to professional education \citep{Beck1999}.

The current situation for children who speak a language other than German at home has changed significantly. These children now start elementary school alongside their monolingually raised peers and, over time, often become more proficient in German than in their heritage language. As a result, the heritage language typically receives input only in the home environment, placing it at a disadvantage.

To help support heritage language development, some schools offer dedicated classes. For example, Turkish heritage language classes are provided by consulates or federal states in every German state where heritage language education is available. However, unlike long-established foreign languages such as Spanish or French, Turkish is rarely taught as a regular foreign language in Germany. Instead, it is more commonly offered as a language of origin in the context of family language classes. At the elementary level, Turkish is typically taught as \textit{Herkunftssprachenunterricht} (`heritage language instruction') or \textit{muttersprachlicher Ergänzungsunterricht} (`additional mother tongue teaching'), continuing until the end of grade 4 (and until the end of grade 6 in states such as Berlin and Brandenburg) \citep{KuppersEtAl2014}.

Only a few states, including North Rhine-Westphalia, Berlin, Bremen, Hamburg, Bavaria, and Lower Saxony, offer Turkish language instruction as a second or third foreign language at the lower and upper secondary levels \citep{KuppersEtAl2014}. Additionally, some states, such as North Rhine-Westphalia, Hamburg, and Bremen, provide Turkish as a foreign language (TaF), however, this option is rarely implemented due to a shortage of trained teachers \citep{Schroeder2003}.

\largerpage
The vitality of the Turkish language in Germany, coupled with the large Turkish population, has led to a significant increase in studies investigating Turkish as spoken in Germany. These studies have explored various linguistic domains of heritage Turkish and found that it often differs from the Turkish spoken in Turkey. For instance, in the domain of phonology, \citet{Queen2001} identified a novel intonation pattern used by Turkish-German bilingual children in both their Turkish and German. The researcher argues that this novel intonation pattern results from the fusion of German and Turkish prosodic features. Similarly, \citet{Kuhn2019} discovered that heritage speakers of Turkish tend to transfer functionally related prosody from German into Turkish, thereby contributing to the contextualization of pragmatic concepts that are not expressed through other linguistic means in Turkish. In morphology, evidence from studies on Turkish heritage speakers in Germany comes primarily from real-time processing experiments. Heritage speakers generally perform similarly to monolingually-raised Turkish speakers in processing derived and inflected morphologically complex words (\citealt{JacobKırkıcı2016}, \citealt{JacobEtAl2019}), regular and irregular aorist verbs (\citealt{UygunClahsen2021}, \citealt{UygunEtAl2023}), grammatical aspect (\citealt{BuyukyildirimEtAl2025}), showing sensitivity to imperfective and perfective distinctions, and grammatical case \citep{OzsoyEtAl2023}. At the same time, they exhibit greater variability with irregular forms, such as irregular aorist verbs (\citealt{UygunClahsen2021}). Comprehension is also influenced by individual factors, with higher language proficiency and faster processing speed associated with higher accuracy rates (\citealt{BuyukyildirimEtAl2025}). Evidence from morphosyntax, by contrast, highlights differences between heritage and majority speakers. For example, heritage speakers in Germany tend to overuse overt subjects compared to monolinguals in Turkey and heritage speakers in the U.S. \citep{OzsoyEtAl2025}. Similarly, \citet{Uygun2022} shows that heritage speakers overuse plural-marked verbs in optional subject-verb agreement contexts. In the domain of syntax, studies have shown that hypotactic means of clause combining are less frequent in Turkish spoken in Germany compared to that in Turkey (\citealt{Treffers-dallerEtAl2006}, \citealt{Bayram2013}, \citealt{IefremenkoSchroeder2019}, \citealt{Schroeder2025}). Conversely, some research indicates a higher frequency of paratactic clause combining in Turkish in Germany compared to majority Turkish (\citealt{Dollnick2013}, \citealt{Guven2016}) and the emergence of novel coordinators in heritage Turkish in Germany (\citealt{Herkenrath2007}, \citealt{Karakoç2007}, \citealt{HerkenrathKarakoc2007}, \citealt{Csato2016}, \citealt{OzsoyEtAl2022}, \citealt{IefremenkoOezsoySchroeder2024}). These differences suggest that heritage speakers may develop new clause combining strategies, likely under the influence of the majority language.

\newpage
In summarizing some of these studies, \citet[34]{KuppersEtAl2015} suggest that the observed changes across different linguistic domains in heritage Turkish in Germany “could potentially be considered characteristics of an emerging new variety: ‘Germany Turkish’.”

\section{Kurmanji as a heritage language in Germany}\label{sec:key:2.2}

An interesting fact of the Turkish migrant population in Germany is that a significant number of these individuals are actually ethnic Kurds. As previously noted, the first wave of immigration began in the 1960s, when primarily male workers from Turkey arrived in Germany as guest workers. The second wave of Kurdish migration to Western Europe occurred in the 1980s, largely consisting of asylum seekers fleeing Turkey \citep{Ostergaardnielsen2001}. Kurdish migration to Germany has continued to the present day, driven by factors such as family reunification, job relocation, and asylum seeking (\citealt{DemmrichArakon2021}, \citealt{Mcgee2018}). Despite their ethnicity, migrants from Turkey are often classified as Turks based solely on their country of origin. Today, there are approximately 1,300,000 Kurds living in Germany (\citealt{GhaderiAlmstadt2023}). However, it remains unclear how many of these individuals identify as Kurds and what language(s) they speak at home.

The situation regarding Kurdish heritage education is complex for several reasons. First, Kurds do not have an officially recognized state, with their population dispersed across countries such as Turkey, Iran, Iraq, and Syria. This raises the issue of which consulate would be responsible for providing Kurdish heritage language education in Germany, particularly since Turkish is the only official language in Turkey. Secondly, Kurdish is a macro-language that encompasses a continuum of languages that are not fully mutually intelligible (\citealt{Hassanpour1992}, \citealt{Sheyholislami2015}). The differences among Kurdish languages are further exacerbated by the fact that they are written in at least three different scripts: the Latin-based alphabet used by Kurds in Turkey (and some in Syria), the Arabic-based alphabet employed by Kurds in Iran and Iraq, and the Cyrillic alphabet used by Kurdish speakers in countries of the former USSR \citep{Sheyholislami2009}. Consequently, it is essential to determine which variant of the Kurdish language should be taught to Kurdish children in Germany. Finally, the lack of trained Kurdish teaching staff poses an additional challenge, as there are currently no departments in German universities that offer teacher training programs specifically for Kurdish educators.

\newpage
Despite these challenges, heritage Kurdish education is offered in several German states. The initiative began in the early 1990s in Lower Saxony, followed by the city-states of Bremen and Hamburg \citep{Schmidt2000}. Currently, Kurdish language education is available in Brandenburg, Bremen, Lower Saxony, North Rhine-Westphalia, and Rhineland-Palatinate (\citealt{Mediendienst2019}). In 2019, the Berlin Senate approved the introduction of Kurdish classes at three schools in Berlin \citep{Derince2022}.

However, the absence of a standardized Kurdish language raises the ongoing issue of which language or dialect to teach. Different federal states have proposed various solutions to address this challenge. For instance, in Bremen, most families whose children participate in Kurdish mother tongue lessons are immigrants from Mardin, Turkey \citep{Schmidt2000}. As a result, researchers and Kurdish language teachers decided to instruct children in the variety most commonly spoken at home. In Lower Saxony, parents were asked to complete a questionnaire regarding their preferred variety for their children's instruction \citep{Schmidt2000}. In Hamburg, although the Kurdish community is quite heterogeneous, Kurdish lessons are organized exclusively for Kurmanji speakers, as the Sorani-speaking population is more dispersed throughout the city \citep{Schmidt2000}.

Probably due to the challenges mentioned earlier, only a limited number of studies have focused on heritage Kurmanji in Germany. One of the few studies examining Kurmanji-Turkish-German trilinguals in Germany, conducted by \citet{Herkenrath2021}, investigates the use of post-predicate elements across the three languages in contact, noting that the employment of the post-predicate position is influenced by various factors: discourse planning and information structure in Turkish, argument structure in Kurmanji, and syntactic structure in German. The researcher observes that the reasons for using the post-predicate position in the languages studied overlap, which can be attributed to their mutual influence through language contact. This exploratory study highlights the need for a more comprehensive investigation of ongoing changes in heritage Kurmanji.

In addition to this research, a small body of studies on Kurmanji in contact with German has addressed issues such as child language acquisition, particularly in preschool-age learners in Germany (\citealt{Kirgiz2022}), literacy acquisition (\citealt{Simsek2016}, \citealt{CorvachoEtAl2021}), language education \citep{Derince2022}, and sociological factors \citep{Eccarius2019}.

Moreover, it is noteworthy that many members of the Kurmanji-speaking population in Germany, whose ancestors immigrated from Turkey, are often able to speak Turkish. For various reasons, Kurmanji-Turkish-speaking parents in Germany may expose their children to both Kurmanji and Turkish. This exposure is frequently attributed to the high vitality of the Turkish language in Germany and its presence in informal public contexts. As a result, children may be raised trilingually, speaking Kurmanji and Turkish at home and in the neighborhood while using German as the majority language in the broader community.

Therefore, when discussing Turkish as an emerging variety in Germany, it is essential to include both Turkish-German bilinguals and Turkish-Kurmanji-German trilinguals in the analysis. To better understand the underlying factors contributing to potential changes, it is crucial to distinguish between these two groups, which will be addressed in the current study.

\section{Kurmanji as a heritage language in Turkey}\label{sec:key:2.3}

As evident from the previous discussion, the contact between Turkish and Kurmanji did not originate in Germany. For centuries, Kurds have lived in a multilingual environment, resulting in the Kurdish language maintaining close contact with various languages, including Arabic, Aramaic, Persian, and Turkish. During the Ottoman Empire, all nations were able to preserve and develop their languages. However, following the establishment of the Republic of Turkey in 1923, significant changes were made to language policy (\citealt{Yagmur2001}). At that time, the Turkish language—specifically the standard variety known as Istanbul Turkish—became the central focus of language policy in the newly formed country. This policy aimed to promote linguistic homogeneity and national unity by establishing Turkish as the sole official language.

The implementation of a single official language led to the marginalization of linguistic diversity in Turkey. As \citet{Yagmur2001} notes, the majority of linguistic minority groups assimilated into the mainstream community. However, due to their significant population and geographic isolation, Kurdish communities managed to maintain their languages (\citealt{Yagmur2001}). \citet{Dorleijn1996} also indicates that in some rural towns in southeastern Turkey, Kurdish served as a \textit{lingua franca} among speakers, including native Turkish speakers, during the 1990s. Consequently, this study emphasizes that Kurmanji can be considered a heritage language, provided that the language is not widely used in public contexts and the majority of the population does not speak Kurmanji.

The situation surrounding the education of Kurmanji in Turkey has undergone significant transformations over the past century. For more than half of the 20th century, there was a concerted campaign aimed at assimilating ethnic minorities into Turkish society, promoting ethnic mixing through both voluntary and forced migration. During this time, the use of Kurmanji faced severe restrictions, it was not only prohibited in education and broadcasting by law but also banned in everyday speech (\citealt{Zeydanlioglu2012}).

However, the 1990s marked a slight shift in this trend, as Kurmanji was no longer classified as a prohibited language. Radio and television broadcasting, as well as the publication of newspapers and books in Kurmanji, were gradually permitted again, albeit with certain restrictions \citep{Derince2013}. Starting in 2003, private courses for teaching Kurmanji were allowed, but these courses were limited to just 10 weeks, with a maximum of 18 hours per week, and were available only for adults (\citealt{Zeydanlioglu2012}).

In 2012, further reforms enabled middle school students (grades 5--8) to choose Kurmanji as an elective course in Turkish public schools (\citealt{Arslan2015}, \citealt{Derince2013}). This course could be taught for only two hours per week and would only be offered if enough students expressed interest. However, in recent years, the teaching of Kurmanji in Turkey has faced some challenges and limitations, which have affected the use and development of the language within the formal education system.

Although the contact between Turkish and Kurmanji has lasted longer than that between Turkish and German, there is a limited number of studies exploring the CLI between the former pair. One notable study by \citet{JohansonTurkyilmaz2006} investigated the influence of Turkish on Kurmanji phonology, finding that the contact of the Sinemili dialect of Kurmanji with Turkish has led to the extensive use of front rounded vowels by speakers of this dialect. Research has also documented instances of lexical borrowing from Turkish into Kurmanji, as shown in studies by \citet{Bulut2006}, \citet{Dorleijn1996, Dorleijn2006}, and \citet{Cabuk2020}. Furthermore, CLI from Turkish on Kurmanji has been observed in the realm of morphology. For instance, \citet{Dorleijn1996} analyzed the speech of two Kurmanji-Turkish bilinguals from Diyarbakır, revealing a significant impact of Turkish on Kurmanji across all linguistic domains, particularly in morphology. This influence was evident through the adoption of Turkish possessive suffixes in Kurmanji, the loss of gender distinctions, and the omission of overt subjects, among other changes. Additionally, \citet{Cabuk2019} noted the borrowing of the Turkish suffix \textit{-mış} into Kurmanji, which loses its original evidential and perfective meanings in the process.

A comprehensive study by \citet{Haig2007} examined the potential influence of Turkish on the Tunceli dialect of Kurmanji across various linguistic domains. The researcher identified instances of lexical borrowing, including discourse markers, conjunctions, and relational nouns. \citet{Haig2007} also suggested that some structural changes in Kurmanji due to Turkish influence are evident in specific constructions, such as those involving the verb \textit{bûn} (‘to be’), used in both static and processual senses. Nevertheless, the author concludes that despite centuries of coexistence between Kurmanji and Turkish speakers, the core grammar of both languages remains distinct \citep{Haig2007}.

\largerpage
Thus, existing studies on the language contact between Kurmanji and Turkish have predominantly focused on the impact of Turkish on Kurmanji, while the influence of Kurmanji on Turkish has received less attention in academic research. One study by \citet{PolatSchallert2013} examined the attainment of a “native-like” Turkish accent among L1 Kurmanji-speaking adolescents in a Turkish-dominant city. The study revealed considerable variation in the accents of the bilingual speakers in their Turkish. The findings indicated that the degree of identification with the Turkish-speaking community positively predicted a “native-like” Turkish accent. Additionally, \citet{Yonat2023} explored the influence of Kurmanji on Turkish in the realm of phonology, uncovering instances of sound transfer from Kurmanji to Turkish. This study showed that sounds characteristic of Kurmanji were incorporated into the Turkish spoken by Kurmanji speakers. The author emphasizes that the Turkish used by Kurmanji speakers differs from Standard Turkish and posits that it represents a distinct native variety within the Turkish language. Furthermore, \citet{Aykac2016} identified several instances of construction copying from Kurmanji into Turkish in the written texts of Kurmanji-speaking students. For example, the author notes that constructions featuring the reflexive pronoun \textit{kendi} (‘own, self’), such as \textit{kendime kitap okuyorum} (literally, ‘I’m reading a book to myself’), are common in the speech of Kurmanji-Turkish bilinguals. In contrast, the same sentence in standard Turkish would typically be expressed without the reflexive pronoun, as in \textit{kitap okuyorum} (‘I’m reading a book’). Aykaç further suggests that such constructions are transferred from Kurmanji, where the reflexive pronoun \textit{xwe} is required in the same context, as demonstrated in the sentence \textit{ez ji xwe re kitêb dixwînim} (‘I’m reading a book’). The author identifies over twenty such constructions in total and argues that they have gradually gained wider currency within the community, to the extent that they are now also used by speakers without a Kurmanji background. From a usage-based perspective, \citet{Aykac2016}'s findings suggest that the adoption of these novel constructions by a larger number of speakers has resulted in greater entrenchment of these forms in the mental representations of individuals, which has, in turn, facilitated the propagation of these constructions throughout the community.

Finally, although there are almost no studies that systematically investigate the societal perception of the Turkish language spoken by bilingual speakers in Turkey, a bias exists within society toward Kurmanji-Turkish bilinguals, particularly regarding the belief that their language does not conform to the norms of Standard Turkish (\citealt{DemirciKleiner1999}). \citet{Tulgar2010}, as noted in \citet{Inal2019}, highlights the ongoing discrimination against Kurds due to the phonetic peculiarities in their Turkish. For example, \citet{Sahin-firat2010} conducted interviews with Kurdish students, some of whom shared experiences of anxiety related to being made fun of because of their pronunciation in Turkish, or the embarrassment they felt when their parents spoke Turkish with an accent at school.

\section{Summary}\label{sec:key:2.4}
\largerpage
This chapter has outlined the sociolinguistic and educational contexts in which Turkish functions as a heritage language in Germany, Kurmanji is a heritage language in Germany and Turkey, and Turkish as a majority language of monolingual and bilingual speakers in Turkey. The overview demonstrates that both languages exist in complex contact situations that often go beyond bilingualism, involving not only two but sometimes three languages in contact, as in the case of Kurmanji-Turkish-German speakers.

The speaker groups under investigation also differ in the intensity of language contact, which can be assessed in terms of both its duration and the size of the communities involved. In Germany, contact between Turkish or Kurmanji and the dominant language German began in the second half of the twentieth century with the arrival of the so-called guest workers. In contrast, contact between Turkish and Kurmanji started before the Ottoman Empire and intensified following the establishment of the Turkish Republic, when Turkish became the country’s sole official language (\citealt{Yagmur2001}). In addition, the Turkish-speaking community in Germany constitutes one of the largest immigrant groups, particularly in Berlin, where the present study’s data were collected. The number of trilingual Kurmanji-Turkish-German speakers, by comparison, is considerably smaller and not precisely known, while the number of Kurdish speakers in Turkey far exceeds that of immigrant communities in Germany. Thus, considering both the duration of contact and community size, the speaker groups in Germany stand in contrast to the Kurmanji-Turkish bilingual group in Turkey investigated in this study.

Finally, the status of the languages under investigation differs across contexts. In Germany, both Turkish and Kurmanji function as heritage languages, in Turkey, Kurmanji can also be considered a heritage, and Turkish serves as the majority language of both monolingual and bilingual speakers. 

Taken together, the sociolinguistic profiles and contact histories discussed in this chapter highlight the diversity of contact settings in which Turkish and Kurmanji are used. Examining CLI across these different groups provides an opportunity to explore how, in addition to linguistic factors, social factors such as language status and the intensity of contact shape the linguistic outcomes of both heritage and majority speakers.

